% Pages liminaires : titre et crédits

\thispagestyle{vide}
\renewcommand{\LJmois}{}\renewcommand{\LJfolio}{}
\begin{center}
\vspace*{14mm}
{\large\scshape Hippolyte Combes}\\[1mm]
{\small (1876–1970)}\\[16mm]
{\fontsize{30pt}{32pt}\selectfont Livre-Journal}\\[5mm]
{\large\itshape 25 juillet 1916 – 30 avril 1917}\\[9mm]
{\normalsize Carnet d'un interprète militaire\\ auprès de l'armée britannique}\\[14mm]
\rule{24mm}{0.4pt}\\[8mm]
{\normalsize Texte diplomatique en regard des fac-similés,\\[1mm]
avec commentaire (critique \& historique),\\[1mm]
cartes et illustrations}\\[16mm]
{\small Édition établie par}\\[1.5mm]
{\large Louis Hauseux}\\
\vfill
{\small 2026}
\end{center}
\clearpage
\thispagestyle{vide}
\vspace*{\fill}
{\footnotesize\setlength{\parindent}{0pt}\setlength{\parskip}{1.2mm}
\textit{Livre-Journal} autographe d'Hippolyte Combes, conservé dans sa descendance : gardé chez son fils François Combes, puis transmis par Maurice Combes à Pierre Combes, petit-fils de l'auteur. Le carnet a été photographié le 26 novembre 2025 (186 clichés) ; ces photographies sont reproduites en regard de la transcription.

Les archives familiales et militaires citées (état signalétique des services, lettres de guerre, portraits, actes d'état civil, biographie par Pierre Combes) sont conservées dans la famille ; leurs reproductions numériques sont réunies, avec les sources de la présente édition, dans le dépôt \url{https://github.com/Ludwig-H/Hippolyte-Combes}.

Format : 160 × 226 mm (proportions A5 ; impression en livret sur feuilles A4 à 92,5 \%). Composition en Libertinus Serif avec XeLaTeX.

\vspace{3mm}
© Louis Hauseux, 2026.}
\clearpage
